% Generated draft: requires mathematical and source review.
\documentclass[11pt]{article}
\usepackage[margin=25mm]{geometry}
\usepackage{fontspec}
\setmainfont{Latin Modern Roman}
\usepackage{amsmath,amssymb,mathtools,unicode-math}
\setmathfont{Latin Modern Math}
\usepackage{enumitem,needspace,xurl,hyperref}
\setlength{\parindent}{0pt}
\setlength{\parskip}{5pt}
\setlength{\emergencystretch}{3em}
\newcommand{\N}{\mathbb{N}}
\newcommand{\Z}{\mathbb{Z}}
\newcommand{\Q}{\mathbb{Q}}
\newcommand{\R}{\mathbb{R}}
\newcommand{\C}{\mathbb{C}}
\newcommand{\F}{\mathbb{F}}
\newcommand{\A}{\mathbb{A}}
\newcommand{\PP}{\mathbb{P}}
\newcommand{\E}{\mathbb{E}}
\newcommand{\eps}{\varepsilon}
\newcommand{\dd}{\,\mathrm{d}}
\begin{document}
\section*{UnsolvedMath ID: 30002992}
\textbf{Source problem code:} catalogue entry
\subsection{Definitions and mathematical statement}
For positive integers \(Q,n\), let
\[
\mathcal A_Q(\R^n)=(\R^n)^Q/\mathfrak S_Q,
\]
where \(\mathfrak S_Q\) acts by permuting coordinates. Thus an element is an unordered collection of \(Q\) points, counted with multiplicity. Its standard matching metric is
\[
\mathcal G([a_1,\ldots,a_Q],[b_1,\ldots,b_Q])
=\min_{\sigma\in\mathfrak S_Q}
\left(\sum_{i=1}^Q |a_i-b_{\sigma(i)}|^2\right)^{1/2}.
\]
A \(Q\)-valued map on a domain \(\Omega\subset\R^m\) is a map \(u:\Omega\to\mathcal A_Q(\R^n)\). These definitions provide terminology; the source does not specify \(m,n,Q\), the domain, or the admissible function space.

For a functional \(\mathcal F\) and a specified competitor class \(\mathcal C(u)\), the minimizing property means
\[
\mathcal F(u)\leq\mathcal F(v)
\qquad\text{for every }v\in\mathcal C(u).
\]
The competitor class encodes the boundary conditions or the permitted local variations, neither of which is supplied here.

\textbf{Source question.} Do minimizers of quasiconvex functionals on \(Q\)-valued functions have partial regularity, and for which subclasses of functionals do regularity results hold?

The requested partial regularity concerns regularity away from an exceptional set. The source does not specify the regularity of the local sheets, the treatment of branching, or the required size of that exceptional set. It also supplies no definition of quasiconvexity for the multivalued setting, no integrand class, and no growth, coercivity, ellipticity, or differentiability hypotheses. Consequently, the supplied question determines a research direction but does not determine a single quantified regularity assertion. A precise theorem or conjecture requires these choices; none is imposed in this draft.
\subsection{Short English statement}
The problem asks whether minimizing maps taking unordered \(Q\)-tuples of values are regular away from a suitably small exceptional set, and which subclasses of quasiconvex functionals admit such results. The functional class, admissible variations, boundary conditions, and intended regularity conclusion remain unspecified.
\subsection{Sources}
\begin{itemize}
\item Local UnsolvedMath registry, supplied record with permanent ID \(30002992\), source code \texttt{catalogue entry}, titled \emph{Partial Regularity of Multivalued Quasiconvex Minimizers}, in the Oberwolfach Reports Open Problems collection. No individual proposer is identified. The supplied proposed year is 2015.
\item Supplied report citation: \emph{Mini-Workshop: Scales in Plasticity (2016)}, Oberwolfach Reports, pp.\ 3007--3048, DOI \texttt{10.4171/owr/2015/53}. \url{https://doi.org/10.4171/owr/2015/53}. The background attributes the question to mathematical workshop and open-problem discussions documented in Oberwolfach Reports.
\item Supplied literature source, labeled ``Oberwolfach Q-valued report'': \url{https://ems.press/content/serial-article-files/46600}.
\item Supplied literature source, labeled ``Dirichlet Q-valued regularity'': \url{https://arxiv.org/abs/0803.0060}.
\item Supplied literature source, labeled ``Q-graph ellipticity framework'': \url{https://arxiv.org/abs/2603.20788}.
\item Literature assessment dated exactly \texttt{2026-08-21}: the supplied curation labels the problem partially solved and states that the general partial-regularity problem remains open, while Dirichlet minimizers and several structured elliptic or polyconvex subclasses have substantial regularity theory. This is supplied dated curation derived from the OWR material, not an independently verified literature assessment. The record requests independent verification of time-sensitive status and cited claims.
\end{itemize}
\end{document}
